\section{Umsetzung in Simulink und Batterieladung}
\label{sec:simulink}

% VERANTWORTLICH: ...
% Aufgabenpunkt 7.

\subsection{Aufbau des Blockschaltbilds}

% Systematik aus der Vorlesung: pro Zustandsvariable ein Integrator,
% davor ein Summenblock, dann die rechte Seite verdrahten.
% Bei nur einem Zustand ist das ein Integrator.
% Auf algebraische Schleifen achten: die Rueckkopplung von T_m auf W_el
% laeuft ueber den Integrator und ist daher unkritisch.

\subsection{Kopplung mit dem Batteriemodell}

\subsection{Vergleich mit der MATLAB-Implementierung}

% Beide Implementierungen mit denselben Eingangsdaten rechnen und die
% Differenz zeigen. Das ist eine echte Verifikation und deutlich mehr wert
% als zwei separate Kurven.
% Keine Simulink-Scopes ins Protokoll. Daten ins Workspace loggen und in
% MATLAB plotten.
